\xchapter{Notes of a Visit to the Russian Church by the Late William Palmer, M.A. Selected and Arranged by Cardinal Newman Kegan Paul \& Co., 1882}

\xsection{Prefatory Notice}

W\textsc{illiam} P\textsc{almer},
Fellow of Magdalen College, Oxford, eldest son of the Rev. William
Jocelyn Palmer, Rector of Mixbury, and brother of Lord Chancellor
Selborne, the Rev. George Horsley Palmer, and Archdeacon Palmer of
Oxford, was one of those earnest-minded and devout men, forty years
since, who, deeply convinced of the great truth that our Lord had
instituted, and still acknowledges and protects, a visible
Church—one, individual, and integral—Catholic, as spread over the
earth, Apostolic as co-eval with the Apostles of Christ, and Holy, as
being the dispenser of his Word and Sacraments—considered it at
present to exist in three main branches, or rather in a triple
presence, the Latin, the Greek, and the Anglican, these three
being one and the same Church, distinguishable from each other only by
secondary, fortuitous, and local, though important, characteristics.
And, whereas the whole Church in its fulness was, as they believed, at
once and severally Anglican, Greek, and Latin, so in turn each one of
those three was the whole Church; whence it followed that, whenever
any one of the three was present, the other two, by the nature of the
case, were absent, and therefore the three could not have direct
relations with each other, as if they were three substantive bodies,
there being no real difference between them except the external
accident of place. Moreover since, as has been said, on a given
territory, there could not be more than one of the three, it followed
that Christians generally, wherever they were, were bound to
recognize, and had a claim to be recognized by, that one, ceasing to
belong to the Anglican Church, as Anglican, when they were at Rome,
and ignoring Rome as Rome, when they found themselves at Moscow.
Lastly, not to acknowledge this inevitable outcome of the initial idea
of the Church, viz., that it was both everywhere and one, was bad
logic, and to act in opposition to it was nothing short of setting up
altar against altar, that is, the hideous sin of schism, and a
sacrilege.

This I conceive to be the formal teaching of
Anglicanism; this is what we held and professed in Oxford forty years
ago; this is what Mr. Palmer intensely believed and energetically
acted on when he went to Russia. It was his motive-cause for going
there; for he hoped to obtain from the Imperial Synod such a
recognition of his right to the Greco-Russian Sacraments, as would be
an irrefragable proof that the doctrine of the Anglican divines was no
mere theory, and that an Anglican Christian was \emph{ipso facto} an
Oriental Orthodox also.

How Mr. Palmer's appeal for such a recognition of
our ``Anglo-Catholicism'' was met by the ecclesiastical authorities of
Petersburg is the main subject of this volume, though not the
main object of its publication. It is published for the vivid picture
it presents to us, for better or for worse, of the Russian Church,
gained, as it was, without effort by the author's intercourse with
priests and laymen, and with the population generally. As might be
expected, they disallowed his claim; but, what was hardly to be
expected, they felt no sympathy for his conception of the Church of
Christ, in its necessary unity, which, even if novel and strange,
could not have been altogether new to them, as being at least part of
that ancient teaching which they so proudly claimed as their own
peculiar prerogative.

Mr. Palmer demanded communion, not as a favour,
but as a right; not as if on his part a gratuitous act, but as his
simple duty; not in order to become a Catholic, but because he was a
Catholic already. Now, if in refusing him they had confined themselves
to the reason which they did also give, that, till he anathematized
the Anglican Articles, they could not be sure he was not a
Lutheran or a Calvinist, they would at least have been intelligible;
or, if they had simply urged, as they also did, that they could not
commit themselves to new precedents for the case of an individual, and
that Synods must meet, and formal correspondence ensue, and
authoritative canons pass, on the part both of Russia and England,
before any acts of communion could take place, that too was a prudent
and sensible course, and would give hopes for the future; but, instead
of keeping to ground so clear and so easily maintained, some of their
highest prelates and officials go out of their way to deny altogether,
or at least to ignore, the Catholicity of the Church as recognized in
the Creed, as if their own time-honoured communion was but a revival
of the ancient Donatists. They say virtually, even if not expressly, ``We
know nothing about Unity, nothing about Catholicity; it is no term of
ours; it had indeed a meaning once, it has none now. Our Church is
not Catholic, it is Holy and Orthodox; also, (because it came from the
East, whence Divine Truth has ever issued,) it is Oriental. We know of
no true Church besides our own. We are the only Church in the world.
The Latins are heretics, or all but heretics; you are worse; we do not
even know your name. There is no true Christianity in the world except
in Russia, Greece, and the Levant; and, as to the Greeks, many as they
are, after all they are a poor lot.''

Let me not be supposed to impute to those
distinguished personages any discourtesy, whether of language or of
conduct, in their intercourse with Mr. Palmer. They gave him a
welcome, which, considering how little they could at first understand
his motives in coming among them, tells altogether in their favour;
they listened to him with interest and earnestness, and, though
political reasons were doubtless on the side of their being courteous
to an Englishman, they were, as if by nature and habit, as frank
and communicative in their conversations with him, as he was on his
part with them. In consequence of this mutual good understanding, Mr.
Palmer made many friends in Russia, and had no reason to regret his
going there. He liked the people and country, and returned there again
and again; and, though he failed from first to last in the direct
object which started him on these expeditions, yet labours such as
his, so Christian in their aim, so disinterested and self-sacrificing
in their circumstances, are, in a religious point of view, never
wasted, never lost. Mr. Palmer's earnest witness to the divine promise
that the Christian Church, unlike the Jewish, should be spread all
over the earth as Catholic and Ecumenical (however defective was his
conception, as an Anglican, of its unity), had from the first its
measure of success in Russia, and that success, whether greater or
less, would of necessity tell upon the theological schools; moreover it would be the more important because it took place at a
time when the so-called Tractarians had, independently of him, been
inculcating the same great truth on their own people in England. It is
no wonder then, that, struck by this coincidence, there were those in
both countries who listened to a preaching which (as far as it
proclaimed the Unity and the Catholicity of the Church,) was as
primitive as it was out of date, and were led on in consequence to
imagine, if not to contemplate, such a union in doctrine and worship
of their respective Churches as would go far towards fulfilling the
idea of a Catholic communion.

I have no temptation, and am in no danger, of
committing myself to extravagant or over-sanguine speculations in such
a matter. Here I agree with Mr. Wallace in his instructive and
interesting work on Russia; a real and effectual union at this time is
a simple \emph{chimera}. ``Of late years,'' he says, ``there has been a
good deal of vague talk about a possible union of the Russian
and Anglican Churches. If by 'union' is meant, simply, union in the
bonds of brotherly love, there can be of course no objection to any
amount of \emph{pia desideria}; but, if anything more real and
practical is intended, I may warn simple-minded, well-meaning people
that the project is an absurdity,'' vol. ii. pp. 194, 195. Of course I
do not sympathize in the tone of this passage; after all, \emph{pia
desideria} are not bad things, though nothing comes of them,—at
least though nothing comes of them at once; however, as to the future,
I am bound to ask all ``men of good will,'' who pray for peace and
unity, whether here or in the North, to ponder the words of a leading
Russian authority introduced into this volume, to the effect that, ``if
England would approach the Russian Church with a view to an
ecclesiastical union, she must do so through the medium of her
legitimate Patriarch, the Bishop of Rome.''

So much on the contents of this volume, which I
have brought together and put into shape, to the best of my power, out
of the materials and according to the evident intentions of Mr.
Palmer, and, I should add, with the valuable assistance of the Rev.
Father Eaglesim of this Oratory. I need hardly say I have no
acquaintance with the Russian language, a condition, if not necessary,
at least desirable, for my present undertaking; but I have been called
to it, as a religious duty, in the following way:—I had often heard
speak of Mr. Palmer's journals of foreign travel at the date when they
were written; and years after, when he was wont to pay me an annual
visit here in the summer or autumn, the only seasons in which the
English climate was possible to him, I used to urge upon him their
publication. But he never gave me any hopes of it, and I ceased to
trouble him on the subject. After a time his spells of serious
indisposition became so frequent, that when we took leave of each
other, it was on my part with the sad feeling that I was bidding
him a last farewell. At length the end came, in 1879, just before I,
in turn, was to have been his guest at Rome; and then I found to my
surprise that, so far from passing over my wish about his journals, he
had by will left me all his papers. This is how he answered my
importunity, showing a loving confidence in me, though involving me in
an anxious responsibility. Of course he did not anticipate that at my
advanced age I could myself do much; but it will be a true
satisfaction to me, if, as I am sanguine enough to expect, this
volume, illustrative of his first visit to Russia, should prove
interesting and useful generally to Christian readers.

I will say one word more:—I cannot disguise
from myself that to common observers, Mr. Palmer was a man difficult
to understand. No casual, nay, no mere acquaintance would have
suspected what keen affections and what energetic enthusiasm
lived under a grave, unimpassioned, and almost formal demeanour. To
unsympathetic or hostile visitors he was careless to defend, or even
to explain, himself or his sayings and doings; and he let such men go
away, indifferent what they might report or think of him. They would
have been surprised to find that what in conversation they might think
a paradox or conceit in him, was, whether a truth or an error, the
deep sentiment and belief of a soul set upon realities and actuated by
a severe conscientiousness. But, whatever might be the criticisms of
those who saw him casually, no one who saw him much could be
insensible to his many and winning virtues; to his simplicity, to his
unselfishness, to his gentleness and patience, to his singular
meekness, to his zeal for the Truth, and his honesty, whether in
seeking or in defending it; and to his calmness and cheerfulness in
pain, perplexity, and disappointment. However, I do not pretend to
draw his character; apart from all personal attributes, he was
to me a true and loyal friend, and his memory is very dear to me.

J. H. N.

B\textsc{irmingham}, \emph{Easter}, 1882.

P.S.—I add a notice of the principal dates of
Mr. Palmer's life, taken from Mr. Bloxam's Register of Members of
Magdalen College.

William Palmer, 1811, July 12th, born.

1823, went to Rugby School.

1820, matriculated at Magdalen College.

1830, University (Chancellor's) Prize for Latin Verse.

1830, First Class in Classics.

1833, University (Chancellor's) Prize for Latin Essay.

1833-30, Tutor at the University of Durham.

1837-39, University (Oxford) Examiner.

1838-43, College (Magdalen) Tutor,

1855, received into the Catholic Church.

1879, April 5th, died at Rome.

While this volume was passing through the press,
I was grieved to read in the public prints a notice of the death of
Sir. Blackmore, whose name occurs so often in it. He had taken a warm
interest in my work, kindly aided me as he only could, and looked
forward to its perusal, when finished, as recalling various pleasant
memories of a valued friend.
